%%%%%%%%%%%%%%%%%%%%%%%%%%%%% Define Article %%%%%%%%%%%%%%%%%%%%%%%%%%%%%%%%%%
\documentclass{article}
%%%%%%%%%%%%%%%%%%%%%%%%%%%%%%%%%%%%%%%%%%%%%%%%%%%%%%%%%%%%%%%%%%%%%%%%%%%%%%%

%%%%%%%%%%%%%%%%%%%%%%%%%%%%% Using Packages %%%%%%%%%%%%%%%%%%%%%%%%%%%%%%%%%%
\usepackage{geometry}
\usepackage{graphicx}
\usepackage{amssymb}
\usepackage{amsmath}
\usepackage{amsthm}
\usepackage{empheq}
\usepackage{mdframed}
\usepackage{booktabs}
\usepackage{lipsum}
\usepackage{graphicx}
\usepackage{color}
\usepackage{psfrag}
\usepackage{pgfplots}
\usepackage{bm}
\usepackage{microtype}
%%%%%%%%%%%%%%%%%%%%%%%%%%%%%%%%%%%%%%%%%%%%%%%%%%%%%%%%%%%%%%%%%%%%%%%%%%%%%%%

% Other Settings

%%%%%%%%%%%%%%%%%%%%%%%%%% Page Setting %%%%%%%%%%%%%%%%%%%%%%%%%%%%%%%%%%%%%%%
\geometry{a4paper}

%%%%%%%%%%%%%%%%%%%%%%%%%% Define some useful colors %%%%%%%%%%%%%%%%%%%%%%%%%%
\definecolor{ocre}{RGB}{243,102,25}
\definecolor{mygray}{RGB}{243,243,244}
\definecolor{deepGreen}{RGB}{26,111,0}
\definecolor{shallowGreen}{RGB}{235,255,255}
\definecolor{deepBlue}{RGB}{61,124,222}
\definecolor{shallowBlue}{RGB}{235,249,255}
%%%%%%%%%%%%%%%%%%%%%%%%%%%%%%%%%%%%%%%%%%%%%%%%%%%%%%%%%%%%%%%%%%%%%%%%%%%%%%%

%%%%%%%%%%%%%%%%%%%%%%%%%% Define an orangebox command %%%%%%%%%%%%%%%%%%%%%%%%
\newcommand\orangebox[1]{\fcolorbox{ocre}{mygray}{\hspace{1em}#1\hspace{1em}}}
%%%%%%%%%%%%%%%%%%%%%%%%%%%%%%%%%%%%%%%%%%%%%%%%%%%%%%%%%%%%%%%%%%%%%%%%%%%%%%%

%%%%%%%%%%%%%%%%%%%%%%%%%%%% English Environments %%%%%%%%%%%%%%%%%%%%%%%%%%%%%
\newtheoremstyle{mytheoremstyle}{3pt}{3pt}{\normalfont}{0cm}{\rmfamily\bfseries}{}{1em}{{\color{black}\thmname{#1}~\thmnumber{#2}}\thmnote{\,--\,#3}}
\newtheoremstyle{myproblemstyle}{3pt}{3pt}{\normalfont}{0cm}{\rmfamily\bfseries}{}{1em}{{\color{black}\thmname{#1}~\thmnumber{#2}}\thmnote{\,--\,#3}}
\theoremstyle{mytheoremstyle}
\newmdtheoremenv[linewidth=1pt,backgroundcolor=shallowGreen,linecolor=deepGreen,leftmargin=0pt,innerleftmargin=20pt,innerrightmargin=20pt,]{theorem}{Theorem}[section]
\theoremstyle{mytheoremstyle}
\newmdtheoremenv[linewidth=1pt,backgroundcolor=shallowBlue,linecolor=deepBlue,leftmargin=0pt,innerleftmargin=20pt,innerrightmargin=20pt,]{definition}{Definition}[section]
\theoremstyle{myproblemstyle}
\newmdtheoremenv[linecolor=black,leftmargin=0pt,innerleftmargin=10pt,innerrightmargin=10pt,]{problem}{Problem}[section]
%%%%%%%%%%%%%%%%%%%%%%%%%%%%%%%%%%%%%%%%%%%%%%%%%%%%%%%%%%%%%%%%%%%%%%%%%%%%%%%

%%%%%%%%%%%%%%%%%%%%%%%%%%%%%%% Plotting Settings %%%%%%%%%%%%%%%%%%%%%%%%%%%%%
\usepgfplotslibrary{colorbrewer}
\pgfplotsset{width=8cm,compat=1.9}
%%%%%%%%%%%%%%%%%%%%%%%%%%%%%%%%%%%%%%%%%%%%%%%%%%%%%%%%%%%%%%%%%%%%%%%%%%%%%%%

%%%%%%%%%%%%%%%%%%%%%%%%%%%%%%% Title & Author %%%%%%%%%%%%%%%%%%%%%%%%%%%%%%%%
\title{Notes for Chapter 2 of ISLP}
\author{Antonio Miguel Natusch Zarco}
%%%%%%%%%%%%%%%%%%%%%%%%%%%%%%%%%%%%%%%%%%%%%%%%%%%%%%%%%%%%%%%%%%%%%%%%%%%%%%%

\begin{document}
\maketitle
\section{What is Statistical Learning?}
At first glance, we're presented
with a Marketing example, where
a series of \textit{budgets}
(namely \textbf{\textcolor{orange}{TV, radio,}}
and \textbf{\textcolor{orange}{Newspaper}})
are examined to see if they have any effect on
\textbf{\textcolor{orange}{sales.}}

Such budgets, in a more abstract, statistical learning approach
are to be considered \textit{input variables}
(which can also be referred to as \textit{predictors, independent variables,
	features} or \textit{variables}), typically denoted by using the
symbol $X$ alongside as many subscripts as needed.

\textbf{\textcolor{orange}{Sales}}, on the other hand,
is what \textbf{depends on} these \textit{features},
so it is often referred to as \textit{response}
or \textit{dependent variable}, typically
represented by $Y$

A relationship between $Y$ and a series of \textit{features}
$X = (X_1, X_2, \dots, X_p)$ is assumed,
and be written in the general form:

\begin{equation}
	Y = f(X) + \epsilon
\end{equation}

, where:

\begin{itemize}
	\item $f$: Fixed but unknown function of $X_1, \dots, X_p$.
	      Represents the \textit{systematic} information that $X$
	      provides about $Y$.
	\item $\epsilon$: Random \textit{error term}, independent of X and
	      has mean zero.
\end{itemize}

\begin{center}
	\textbf{Statistical learning is all about estimating this}
	$f$ \textbf{function, whether for \textit{prediction}
		or \textit{inference}}.
\end{center}

It is worth noting that $f$, when plotted,
is a $p$-dimensional surface that is estimated
based on the observed data,
where $p$ refers to the amount of \textit{features}
that our studied context possesses.

Also, \textbf{$\epsilon$} averages to zero due
to the following reason:

Defining the sample regression equation for each observation $i$:

\begin{equation}
	y_i = \hat{\beta}_0 + \hat{\beta}_1 x_i + \hat{u}_i
	\label{eq:sample-regression-equation}
\end{equation}

, where:

\begin{itemize}
	\item $y_i$: Response of observation $i$.
	\item $\hat{\beta_0}$: Bias term, also known as the \textit{intercept}.
	\item $\hat{\beta_1}$: Regression coefficient, which measures
	      the expected change of a response $Y$ for a one-unit
	      change in the observation $X$, assuming other variables
	      are constant. In simple linear regression ($Y = \beta X + a$),
	      $\beta$ is the slope of the line of best fit.
	\item $x_i$: Observation being studied, assuming that $p = 1$.
	\item $\hat{u_i}$: Residual or \textit{sample error term}.
\end{itemize}

Rearranging \eqref{eq:sample-regression-equation}, we get
the formula for a single residual:

\begin{equation}
	\hat{u}_i = y_i - \hat{\beta}_0 - \hat{\beta}_1 x_i
	\label{eq:rearrange-sample-reg-eq}
\end{equation}

And, since Ordinary Least Squares works by minimizing
the Sum of Squared Residuals (SSR),
where a single squared residual can
be seen on \eqref{eq:rearrange-sample-reg-eq}:

\begin{equation}
	\textit{SSR} = \sum_{i=1}^{n} \hat{u}_i ^ {2} = \sum_{i=1}^{n} (y_i - \hat{\beta}_0 - \hat{\beta}_1 x_i)^2
\end{equation}

, we can see that minimizing this sum
by taking its partial derivative
through the chain rule:

\begin{equation}
	\frac{\partial\text{SSR}}{\partial \hat{\beta}_0} = \frac{\partial}{\partial\hat{\beta}_0} \sum_{i=1}^{n} (y_i - \hat{\beta}_0 - \hat{\beta}_1 x_i)^2 = 0
\end{equation}

\begin{equation}
	\sum_{i=1}^{n} 2(y_i - \hat{\beta}_0 - \hat{\beta}_1 x_i) \cdot (-1) = 0
\end{equation}

($\cdot -2$):

\begin{equation}
	\sum_{i=1}^{n}(y_i - \hat{\beta}_0 - \hat{\beta}_1 x_i) = 0
	\label{ssr-diff-third-step}
\end{equation}

and since $\hat{u}_i$ is equal to the expression
inside the parentheses on \eqref{ssr-diff-third-step}
as seen on \eqref{eq:rearrange-sample-reg-eq}
$\epsilon$ is indeed 0:

\begin{equation}
	\begin{split}
		\sum_{i=1}^{n}\hat{u}_i                             & = 0 \\
		\therefore \epsilon \equiv  \hat{u}_i \footnotemark & = 0
	\end{split}
\end{equation}
\footnotetext{This is in the context of linear regression.}



% section What is Statistical Learning? (end)
\end{document}
